% Generated draft: requires mathematical and source review.
\documentclass[11pt]{article}
\usepackage[margin=25mm]{geometry}
\usepackage{fontspec}
\setmainfont{Latin Modern Roman}
\usepackage{amsmath,amssymb,mathtools,unicode-math}
\setmathfont{Latin Modern Math}
\usepackage{enumitem,needspace,xurl,hyperref}
\setlength{\parindent}{0pt}
\setlength{\parskip}{5pt}
\setlength{\emergencystretch}{3em}
\newcommand{\N}{\mathbb{N}}
\newcommand{\Z}{\mathbb{Z}}
\newcommand{\Q}{\mathbb{Q}}
\newcommand{\R}{\mathbb{R}}
\newcommand{\C}{\mathbb{C}}
\newcommand{\F}{\mathbb{F}}
\newcommand{\A}{\mathbb{A}}
\newcommand{\PP}{\mathbb{P}}
\newcommand{\E}{\mathbb{E}}
\newcommand{\eps}{\varepsilon}
\newcommand{\dd}{\,\mathrm{d}}
\begin{document}
\section*{UnsolvedMath ID: 700272}
\textbf{Source problem code:} catalogue entry
\subsection{Definitions and mathematical statement}
For distinct prime numbers \(p\) and \(q\), define the positive integer
\[
S(p,q)=\frac{p^q-1}{p-1}=\sum_{i=0}^{q-1}p^i.
\]
Here a prime is a positive integer greater than \(1\) whose only positive divisors are \(1\) and itself. The conjecture proposed by J. G. Thompson is
\[
\forall p,q\text{ prime},\qquad p\ne q\ \Longrightarrow\ S(p,q)\nmid S(q,p).
\]
Thus the assertion concerns every ordered pair of distinct primes, including pairs containing \(2\). Equivalently, there is no integer \(k\geq 1\) such that
\[
\frac{q^p-1}{q-1}=k\frac{p^q-1}{p-1}.
\]
The supplied source states that validity of this conjecture would simplify the proof of solvability of groups of odd order by making detailed use of generators and relations unnecessary. This is preserved as source-provided context, not as an independently verified implication. The registry labels the conjecture open; that label is supplied curation, not an independently verified assessment.
\subsection{Short English statement}
Thompson conjectures that, for every pair of distinct primes \(p,q\), the integer \(1+p+\cdots+p^{q-1}\) does not divide \(1+q+\cdots+q^{p-1}\). The supplied record labels the conjecture open and describes a possible simplification of the proof of solvability of groups of odd order if it holds.
\subsection{Sources}
\begin{itemize}
\item Local UnsolvedMath registry, supplied complete source record with permanent ID \(700272\), exact source code \texttt{catalogue entry}, title \emph{Kourovka Notebook Problem 4.65}. The record attributes the conjecture to J. G. Thompson, proposed in 1973, and supplies the status \texttt{open}.
\item \emph{The Kourovka Notebook (Unsolved Problems in Group Theory)}, Problem 4.65, originating in the 4th issue (1973); the supplied background identifies the 21st-issue edition (2026), \texttt{arXiv:1401.0300v48}. The supplied collection description names E. I. Khukhro and V. D. Mazurov as editors.
\item W. Feit and J. G. Thompson, \emph{Pacific J. Math.}, volume 13, no. 3 (1963), pp. 775--1029. This reference is supplied for the contextual discussion of the proof of solvability of groups of odd order; no article title is supplied.
\end{itemize}
\end{document}
